\documentclass[12pt]{article}
\usepackage[utf8]{inputenc}
\usepackage[T1]{fontenc}
\usepackage{hyperref}
\title{Unubahe Language Resources for K-12 Curriculum Development}
\author{Compiled by Claude Sonnet 5 (claude-sonnet-5) for Andr\'as Kornai}
\date{}
\begin{document}
\maketitle
\section{Introduction}
Unubahe (ISO 639-3 unu) is a nearly extinct Oceanic language at the southeastern tip of Milne Bay Province, Papua New Guinea (Suauic subgroup, Papuan Tip). A 2001 survey counted 7 fluent speakers and 59 with some ability. Almost nothing is published.
\section{Primary Educational Resources}
None found.
\section{Linguistic and Reference Materials}
The main documentation is an unpublished manuscript by Michael J. Rueck, ``Unubahe: A nearly extinct language'' (11 pp., undated), which is cited in Glottolog. \url{https://glottolog.org/resource/reference/id/313374} The Glottolog language page is \url{https://glottolog.org/resource/languoid/id/unub1234.bigmap.html}. ASJP has a wordlist compiled by Viktoria Smirnova. \url{https://asjp.clld.org/languages/UNUBAHE} General description: \url{https://en.wikipedia.org/wiki/Unubahe_language}
\section{Religious and Community Resources}
None found.
\section{Digital and Audio Resources}
None found beyond the ASJP wordlist.
\section{Curriculum Development Context}
There is no basis yet for a curriculum. Documentation with the remaining speakers would have to come first, in consultation with the community.
\end{document}